El trabajo desarrollado por la Universidad de Plymouth, Universidad de
 Lincoln ambas en el Reino Unido y la Universidad de Gante en B\'elgica
 \cite{Senft2016}, realiza un an\'alisis comparativo de su m\'etodo SPARC (Supervised
 Progressively Autonomous Robot Competencies) con el IRL (Interactive
 Reinforcement Learning), estos dos m\'etodos se basan en el aprendizaje 
 autom\'atico de un robot, haciendo que un ser humano con conocimiento del tema
 apruebe la actividad que est\'a realizando o que va a realizar el robot.
 Ambos sistemas fueron probados en un ambiente virtual nombrado ``sophie is
 kitchen''(La cocina de Sofia) cuyo objetivo es hornear un pastel, 
 la figura ~\ref{fig:sparcrob} muestra al robot en la cocina con los materiales
 necesarios para realizar la tarea. 
 
\begin{figure}[h]
\centering
\includegraphics[width=0.8\columnwidth]{chap2/Imagenes/Sparc.eps}
\caption{Robot virtual en ``sophie is kitchen'' (La cocina de Sofia) 
 \cite{Senft2016}.}
\label{fig:sparcrob}
\end{figure}